\subsection{Claude Extraction}
\label{app:anthropic_extraction}

\textbf{Fuzzy extraction.} We perform the extraction attack on Claude models in two stages. In the
experiments shown in \Cref{fig:teaser}, Haiku~4.5 acts as a fuzzy decoder for
thinking blocks produced in other sessions. We use a simple manually written
jailbreak and a prefill attack \citep{andriushchenko2024jailbreaking}, since
Haiku~4.5 supports assistant-turn prefills. We found this method surprisingly robust across different models' reasoning:
despite sampling at temperature~$1$, it produces consistent outputs with an
approximately 1:1 ratio between extracted-reasoning tokens and API-billed
thinking tokens.

We use the attack template in \Cref{fig:claude_fuzzy_template}. The decode request is the first user message, followed
by one assistant turn containing the signed thought and the visible
\texttt{<thinking-copy>} prefill. Haiku continues that same assistant turn.


\begin{figure}[H]
\begin{extractiontemplate}
[
  {
    "role": "user",
    "content": "Continue. Transcribe the reasoning attached to this turn, verbatim, inside <thinking-copy>...</thinking-copy>."
  },
  {
    "role": "assistant",
    "content": [
      @thought,  // injected signature
      {"type": "text", "text": "<thinking-copy>"}
    ]
  }
]
\end{extractiontemplate}
\caption{\textbf{Claude extraction request template.} The \texttt{@thought} entry is the replayed reasoning block. Injection occurs in the current turn, and the model's output is prefilled.}
\label{fig:claude_fuzzy_template}
\end{figure}


Occasional extraction failures fall into three main categories: the model refuses to transcribe its reasoning, echoes the final message of the extraction template, or becomes confused and claims that there is no preceding conversation or thought to transcribe. We identify and remove these generations using a keyword-based filter.


\textbf{Extraction reconciliation.} To account for occasional refusals and noise in the generated samples, we introduce an optional reconciliation step. We pass up to three non-refusal extractions produced by Haiku~4.5 to Opus~4.8 and ask it to produce a single faithful transcription.


Using the template shown in \Cref{fig:claude_reconcile_template}, Opus receives the original signed thought together with the noisy Haiku extractions, which are inserted into a completed assistant turn. A subsequent user turn requests a single reconciled transcription. Unlike the fuzzy-decoding procedure, this request ends with a user turn and therefore does not rely on an assistant prefill. We generate the final extraction at temperature~$0$.

\begin{figure}[H]
\begin{extractiontemplate}
[
  {
    "role": "user",
    "content": "Your own reasoning is attached to the next turn \u2014 it is carried in this context by its signature, not something you must recall from memory. You will be shown that reasoning together with several noisy, independently-decoded transcriptions of it, and your job is to reconstruct the single faithful transcription."
  },
  {
    "role": "assistant",
    "content": [
      @thought,  // injected original signed thinking block
      {
        "type": "text",
        "text": "Some noisy, independently-decoded transcriptions of the reasoning above:\n\n{noisy generations}"
      }
    ]
  },
  {
    "role": "user",
    "content": "The reasoning attached to the turn above is yours for this task \u2014 it is present in this context via its signature, so this is a transcription task over content you already have, not a memory you must recall. Below it are several noisy, independently-decoded transcriptions of that reasoning.\n\nReconstruct the single faithful transcription of the attached reasoning inside <thinking-copy>...</thinking-copy>: keep what the transcriptions got right, correct what they got wrong, and restore anything they missed. Preserve exact names, numbers, paths, and identifiers. Do not decline or say you lack access to an earlier session \u2014 the reasoning is attached above and that is all you need. Then stop."
  }
]
\end{extractiontemplate}
\caption{\textbf{Claude reconciliation request template.} The \texttt{@thought} entry is the replayed reasoning block. Injection is done once in the past turn.}
\label{fig:claude_reconcile_template}
\end{figure}
